
El objetivo de este capítulo es demostrar que toda variedad que admita una función de Morse cuyos conjuntos subnivel sean compactos tiene el tipo de homotopía de un CW-complejo (\ref{cwvariedad}). Con ese fin, introducimos formalmente qué entendemos por CW-complejo y por tipo de homotopía.

Incluimos también dos resultados necesarios para nuestro objetivo, pero interesantes en sí mismos, que describen bajo qué condiciones y cómo cambia el tipo de homotopía a medida que recorremos los conjuntos subnivel de una función de Morse (\ref{retractodef} y \ref{pegadocelula}). 

El capítulo concluye con dos ejemplos de cómo se pueden aplicar los principales teoremas demostrados.


\section{CW-Complejos y tipo de homotopía}
%tipo de homotopia-----------------------
\begin{definition}
Sea $X$ un espacio topológico y $A \subset X$ un subespacio. Decimos que $A$ es un
\textit{retracto de deformación fuerte} de $X$ si existe una homotopía relativa a $A$ de la función identidad de $X$, $id_X$, a una retracción $r: X \to A$.
\end{definition}

\begin{nota}
    A la hora de denotar o construir homotopías usaremos indistintamente dos notaciones. Dados dos espacios topológicos $X$ e $Y$, nos referimos a una homotopía entre dos funciones continuas $f:X\to Y$ y $g:X\to Y$ como la función continua $H(x,t): X\times [0,1]\to Y$ tal que $H(x,0)=f(x)$ y $H(x,1)=g(x)$ para todo $x\in X$ o como la familia de funciones continuas $\{H_t: X\to Y\}_{t\in [0,1]}$ tales que $H_t(x)=H(t,x)$ para cualesquiera $t$ y $x$.
\end{nota}

\bigskip

\begin{definition}
Sean $X$ e $Y$ espacios topológicos. Decimos que $X$ e $Y$ son
\textit{homotópicamente equivalentes} si existen aplicaciones continuas
\begin{equation*}
    f: X \to Y, \qquad g: Y \to X
\end{equation*}
tales que
\begin{equation*}
    g \circ f \simeq \mathrm{id}_X \qquad \text{y} \qquad f \circ g \simeq \mathrm{id}_Y.
\end{equation*}
En este caso, decimos que $X$ e $Y$ tienen el mismo \textit{tipo de homotopía} y llamamos a $f$ y $g$ \textit{equivalencias homotópicas}.
\end{definition}

\medskip

\begin{nota}
    En lo que sigue nos referiremos a los retractos de deformación fuerte simplemente como \textit{retractos de deformación}.
\end{nota}

\begin{remark}
    Nótese que todo retracto de deformación de un espacio tiene el mismo tipo de homotopía que este. Dados un espacio $X$ y un retracto de deformación $A\subseteq X$, la función inclusión, $i: A\hookrightarrow X$, y la retración $r: X \to A$ son equivalencias homotópicas.
\end{remark}
%--------------------------------------------------------

\bigskip

%CW-complejos-------------------------------
\begin{definition}
    Una \textit{celda} $e^n$ es un espacio topológico homeomorfo al disco unidad n-dimensional $D^n$. Diremos que la dimensión de la celda es $n$ y que $e^n$ es una \textit{n-celda}.
\end{definition}

\bigskip

\begin{definition}
    Llamamos \textit{CW-complejo} a un espacio topológico construido siguiendo los siguientes pasos:

        \begin{enumerate}[label=(\roman*)]
            \item Se comienza con un conjunto discreto $ X_0 $ (llamado también \textit{0-esqueleto}), cuyos elementos llamaremos \textit{0-celdas}.
            
            \item Se construye, de manera inductiva, el llamado $ n $-esqueleto para $n>0$ como
            \begin{equation*}
                X_n = X_{n-1}\cup_{\varphi_{\alpha}}D^n_{\alpha},~\alpha\in \Lambda_n
            \end{equation*}
            Es decir, $X_n$ se obtiene de $X_{n-1}$ pegando sobre este una colección indexada por $\Lambda_n$ de $n$-celdas mediante las \textit{funciones de pegado} $\varphi_{\alpha} : \partial D^n_{\alpha} \to X_{n-1}$.
            
            \item Puede pararse tras un número finito de pasos (aunque no tiene por qué hacerse). Llamaremos \textit{CW-complejo} al espacio 
            \begin{equation*}
                X = \bigcup_{n} X_n.    
            \end{equation*}
            Si el número de pasos es infinito, los conjuntos abiertos de la topología serán aquellos que sean abiertos para cada $ X_n $.
        \end{enumerate} 
\end{definition}

\bigskip

\begin{remark} \hfill
    \begin{enumerate}
        \item Decimos que la dimensión de un CW-complejo $X$ es $m$ si la construcción tiene una cantidad finita $m$ de pasos, en ese caso, $X=X_m$. Como se deja entrever en el último paso de la construcción, un CW-complejo puede tener dimensión infinita.
        \item Más formalmente, llamamos $n$-celda a la imagen de la inclusión $i_{\alpha}: D^n_{\alpha}\to X_n$ y dotamos a $X$ de la \textit{topología débil}, es decir, $\mathcal{U}\subseteq X$ es abierto si y solo si $i^{-1}_{\alpha}(\mathcal{U})\subseteq D^k_{\alpha}$ es abierto para todos $\alpha$ y $k\ge 0$. De hecho, es esta topología la que les da nombre, ya que CW viene del inglés ``Closure finite-Weak topology''
    \end{enumerate}
\end{remark}

Los CW-complejos presentan una estructura cómoda para aplicar técnicas algebraicas y computar invariantes topológicos como el grupo fundamental o los grupos de homología, que veremos más adelante. Existe también la ventaja de que toda variedad diferenciable puede dotarse de una estructura de este tipo bien vía triangulaciones (\textit{cf.} sección \ref{subsec:triangulaciones}), o bien vía teoría de Morse.

\begin{nota}
    Llamaremos habitualmente \textit{vértices} a las $0$-celdas, \textit{aristas} a las $1$-celdas y \textit{caras} a las $2$-celdas. Nótese también que el borde o la frontera de una arista son dos vértices y de una cara una colección finita de aristas.
\end{nota}

La seccion \ref{subsec:homcelular} sobre homología celular contiene algunos ejemplos de CW-complejos.

\medskip

\begin{definition}
    A cada celda $e^n_{\alpha}$ de un CW-complejo $X$ le corresponde una \textit{función característica} $\phi_{\alpha} : D^n_{\alpha}\to X$ que puede ser factorizada como
    \begin{equation*}
        D^n_{\alpha} \hookrightarrow X_{n-1} \sqcup_{\alpha} D^n_{\alpha} \to X_n \hookrightarrow X
    \end{equation*}
    donde la función del medio es la función cociente que define $X_n$ (es decir, la que identifica los puntos de $D^n_{\alpha}$ con sus imágenes por $\varphi_{\alpha}$) y las otras inclusiones naturales. $\phi_{\alpha}$ extiende a la función de pegado $\varphi_{\alpha}$ y es un homeomorfismo del interior de $D^n_{\alpha}$ sobre $e^n_{\alpha}$.

\end{definition}
%---------------------------------------------------------------
\bigskip

\begin{proposition}
    Dados el $n$-esqueleto $X_n$ y el $(n-1)$-esqueleto $X_{n-1}$ de un CW-complejo $X$ se tiene que 
    \begin{equation*}
        X_n/X_{n-1} = \bigvee_{\alpha\in\Lambda_n} D^n_{\alpha}/\partial D^n_{\alpha} = \bigvee_{\alpha\in\Lambda_n} \mathbb{S}^{n-1}_{\alpha},
    \end{equation*}
    donde $\vee$ denota la suma wedge \cite[p.~10]{hatcher2002} y $\Lambda_n$ recorre las $n$-celdas de $X$.
\end{proposition}

Esto último se sigue de que cada $n$-celda $e_{\alpha}^n$ se obtiene pegando la frontera de un $n$-disco al $(n-1)$-esqueleto. Al cocientar $X_n/X_{n-1}$ estamos colapsando a un solo punto el $(n-1)$-esqueleto, de manera que componiendo la función de pegado con la de paso al cociente lo que se obtiene es la clase trivial. Así, la frontera del disco queda identificada a un solo punto, además, sabemos que $D^n/\partial D^n \cong \mathbb{S}^n$.

Como las fronteras de todos los  $n$-discos quedan identificadas al mismo punto lo que se obtiene es un ``ramillete'' de esferas, que denotamos
\begin{equation*}
    \bigvee_{\alpha\in\Lambda_n} \mathbb{S}^{n-1}_{\alpha}.
\end{equation*}